% Generated draft: requires mathematical and source review.
\documentclass[11pt]{article}
\usepackage[margin=25mm]{geometry}
\usepackage{fontspec}
\setmainfont{Latin Modern Roman}
\usepackage{amsmath,amssymb,mathtools,unicode-math}
\setmathfont{Latin Modern Math}
\usepackage{enumitem,needspace,xurl,hyperref}
\setlength{\parindent}{0pt}
\setlength{\parskip}{5pt}
\setlength{\emergencystretch}{3em}
\newcommand{\N}{\mathbb{N}}
\newcommand{\Z}{\mathbb{Z}}
\newcommand{\Q}{\mathbb{Q}}
\newcommand{\R}{\mathbb{R}}
\newcommand{\C}{\mathbb{C}}
\newcommand{\F}{\mathbb{F}}
\newcommand{\A}{\mathbb{A}}
\newcommand{\PP}{\mathbb{P}}
\newcommand{\E}{\mathbb{E}}
\newcommand{\eps}{\varepsilon}
\newcommand{\dd}{\,\mathrm{d}}
\begin{document}
\section*{UnsolvedMath ID: 30004741}
\textbf{Source problem code:} catalogue entry
\subsection{Definitions and mathematical statement}
The supplied record poses two conjectural questions for every smooth planar domain \(\Omega\subset\R^2\). Here a domain means a nonempty connected open set. Write \(|\Omega|\) for its area and \(\ell(\partial\Omega)\) for its boundary length. The record does not specify boundedness, the precise boundary regularity meant by smoothness, the spectral operator, its boundary conditions, or the coefficient \(C\). Consequently, the spectral question below remains a formulation with unresolved parameters.

For a specified nonnegative self-adjoint spectral problem on \(\Omega\) with discrete eigenvalues \(\lambda_j\), listed with multiplicity, define
\[
N(\lambda)=\#\{j:\lambda_j\leq\lambda\}.
\]
This definition requires a spectral realization for which the count is finite for each finite \(\lambda\); such a realization is not supplied. In particular, the record does not identify whether the intended problem is the Dirichlet or Neumann Laplacian.

\begin{enumerate}
\item For every smooth planar domain \(\Omega\), is the intended eigenvalue counting function asymptotic to
\[
N(\lambda)=\frac{|\Omega|}{4\pi}\lambda
+C\,\ell(\partial\Omega)\lambda^{1/2}
+o(\lambda^{1/2})
\qquad(\lambda\to+\infty)?
\]
For each fixed admissible domain and spectral realization, the remainder notation means that the difference between \(N(\lambda)\) and the two displayed terms, divided by \(\lambda^{1/2}\), tends to zero. The source leaves the value and permitted dependence of \(C\) unspecified. The expression also requires finite geometric coefficients and an appropriate discrete spectrum, neither of which follows from the explicitly stated smoothness assumption alone.

\item For every smooth planar billiard domain \(\Omega\), does the set of periodic billiard trajectories have measure zero?

To describe a conventional precise interpretation, let
\[
S^1=\{v\in\R^2:|v|=1\},\qquad X_\Omega=\Omega\times S^1.
\]
A regular unit-speed billiard trajectory moves along straight lines in the interior and reflects at a transverse boundary collision according to
\[
v_{\mathrm{out}}=v_{\mathrm{in}}-2\langle v_{\mathrm{in}},n(x)\rangle n(x),
\]
where \(n(x)\) is the outward unit normal. On initial states for which the billiard flow \(\Phi^t\) is defined through the required collisions, set
\[
P_\Omega=\{z\in X_\Omega:\text{there exists }T>0\text{ with }\Phi^T(z)=z\}.
\]
With the phase-space measure \(\mathrm d\mu_\Omega=\mathrm dx_1\,\mathrm dx_2\,\mathrm d\theta\), where \(v=(\cos\theta,\sin\theta)\), the conventional question is
\[
\mu_\Omega(P_\Omega)=0.
\]
The source does not specify this phase space or measure, or its convention for grazing and other singular trajectories. Thus this is an explicitly conditional interpretation of its measure-zero question, requiring confirmation against the report.
\end{enumerate}

The displayed source statements quantify over all smooth planar domains and do not impose convexity. The supplied literature assessment instead singles out the conjecture for every smooth convex table as remaining open; this narrower scope does not replace the quantifiers in the statements. The two questions are presented separately, without asserting their equivalence.
\subsection{Short English statement}
The record asks whether every smooth planar domain has the displayed two-term eigenvalue asymptotic and whether its periodic billiard trajectories have measure zero. The spectral operator, boundary conditions, coefficient \(C\), admissible domain class, and trajectory measure require clarification. The supplied status is partially solved; this draft does not independently verify that assessment.
\subsection{Sources}
\begin{itemize}
\item Local registry: \texttt{data/unsolvedmath/problems.json}, supplied record with permanent UnsolvedMath ID \(30004741\), source code \texttt{catalogue entry}, title \emph{Two-Term Weyl Asymptotics and Ivrii's Billiard Conjecture}, in the Oberwolfach Reports Open Problems collection. No individual proposer is supplied; the proposed year is recorded as 2021.
\item Supplied report citation: \emph{Dynamische Systeme (2022)}, Oberwolfach Reports, pp.\ 1735--1803, DOI \texttt{10.4171/owr/2021/33}, \url{https://doi.org/10.4171/owr/2021/33}. The supplied background describes these reports as documenting mathematical workshops and their open-problem discussions. The citation's 2022 label and its DOI's 2021 component are retained as supplied.
\item Supplied literature reference, labelled ``current billiard Weyl-law results'': \url{https://arxiv.org/abs/2406.16172}.
\item Supplied literature reference, labelled ``2025 algebraic billiards paper'': \url{https://www.sciencedirect.com/science/article/abs/pii/S0001870825005729}.
\item Literature assessment dated exactly \texttt{2026-08-22}: supplied dated curation, not independently verified. It labels the problem partially solved and states that Weyl's two-term formula follows under Ivrii's measure-zero hypothesis and is known in important analytic or generic classes, while Ivrii's conjecture for every smooth convex table remains open. The background characterizes this assessment as dated literature triage derived from the supplied OWR curation and requests independent verification of time-sensitive status and cited claims.
\item The registry labels the statement \texttt{corrected\_verified} and reports that it was checked against the official Oberwolfach report and preserves its hypotheses, notation, and open target. This verification claim is supplied provenance; it has not been independently checked in this draft.
\end{itemize}
\end{document}
